%% =====================================================================
%%  T-06-01  Protest Too Much
%%  -------------------------------------------------------------------
%%  One idea per file. Core is mandatory. Slots in canonical order.
%%  Max 700 words. No layout commands. No filler. New source material
%%  for this entry goes in sources/<BOOK>/T-06-01.tex -- never by
%%  rewriting this file (docs/RULES.md R0.1).
%% =====================================================================
%% @kind: topic
%% @id: T-06-01
%% @title: Protest Too Much
%% @subtitle: Insistence is evidence
%% @chapter: c06
%% @order: 10
%% @status: stable
%% @sources: B1
%% @updated: 2026-09-19

\topic{T-06-01}{Protest Too Much}{Insistence is evidence}

\begin{core}
A quality that has to be asserted repeatedly is usually being asserted because it is absent. Honest people do not advertise honesty, solvent people do not mention solvency, and a person who keeps returning to their own reliability is managing a doubt they know is there.
\end{core}
\begin{mechanism}
Two things produce over-assertion, and they are hard to separate. The first is discomfort: a person who knows a statement is false over-compensates, because the ordinary level of emphasis feels insufficient to carry it. The second is selling: a person proposing an arrangement in which the terms do not demonstrate their good faith must supply the demonstration verbally. In both cases the volume of the claim is doing work that evidence would otherwise do. The tell is therefore not the claim itself but the mismatch between the claim's intensity and the available support --- favourable terms, a verifiable history, a written agreement.
\end{mechanism}
\begin{tells}
  \item Trustworthiness is asserted before it is questioned.
  \item The same assurance is repeated in nearly identical words.
  \item Gratitude and warmth are expressed at length in advance of any benefit received.
  \item Your gain is restated several times by the person proposing the arrangement.
  \item The assurance is emotional rather than factual: tone carries it, not evidence.
  \item Physical warmth increases at the point where the commitment is requested.
\end{tells}
\begin{moves}
  \item Supply the assurance yourself, unprompted and repeatedly, where the terms do not demonstrate good faith.
  \item Attach the assurance to a moment of request rather than to a moment of information.
  \item Restate the other party's benefit until it substitutes for evidence of it.
  \item Increase warmth as the commitment approaches.
\end{moves}
\begin{counters}
  \item Do not argue with the assertion. Ask for the evidence it is replacing: references, a written term, a verifiable history, a small first transaction.
  \item Separate the claim from the deal. Evaluate the terms as though the assurance had never been made.
  \item Treat repeated assurance as a request to stop checking, and check more.
  \item Delay. Over-assertion is time-sensitive; a person selling certainty degrades visibly when asked to wait.
  \item Notice your own over-assertion. If you find yourself insisting on your good faith, you have already lost the position you are defending.
\end{counters}
\begin{cost}
The tell has a high false-positive rate in two populations: anxious people, who over-explain generally, and people who have been badly treated before, who pre-emptively defend themselves. Neither is lying. It is also useless against the skilled operator, who supplies exactly the ordinary amount of emphasis and lets the structure of the deal do the persuading --- which is why this test must be combined with T-05-02 rather than used alone.
\end{cost}

\sources{B1}
\seesources
\seealso{T-03-02, T-05-02, T-06-02}
